\documentclass[11pt]{amsart}
\usepackage[T1]{fontenc}
\usepackage{lmodern}
\usepackage{amsmath,amssymb,amsthm,mathtools}
\usepackage{booktabs}
\usepackage{xcolor}
\usepackage[colorlinks=true,linkcolor=blue!55!black,citecolor=blue!55!black,urlcolor=blue!55!black]{hyperref}
\emergencystretch=1em
\newtheorem{theorem}{Theorem}[section]
\newtheorem{proposition}[theorem]{Proposition}
\title[Eight-Vertex Whiskered-Graph WLP]{The Whiskered-Graph Weak Lefschetz Conjecture\\for Graphs on Eight Vertices}
\author{Zijian Zeng}
\address{Institute of Computer Science and Digital Innovation, UCSI University, Kuala Lumpur, 56000, MALAYSIA}
\email{zijianzeng@foxmail.com; 1002266693@ucsiuniversity.edu.my}
\date{August 31, 2026}
\hypersetup{pdftitle={The Whiskered-Graph Weak Lefschetz Conjecture for Graphs on Eight Vertices},pdfauthor={Zijian Zeng}}
\begin{document}
\begin{abstract}
Over fields of characteristic zero, we verify the eight-vertex case of the
whiskered-graph conjecture of Holleben and Nicklasson. A structural argument reduces
the new computation to graphs of independence number three. There are
6,021 such isomorphism classes. For every class, an exhaustive canonical
enumeration and exact integer arithmetic produce a nonzero kernel vector for
the multiplication map from degree four to degree five. The complete list of
kernel vectors is independently checkable by a short verifier that reconstructs
the multiplication matrices from the graph6 representatives.
\end{abstract}
\maketitle
\section{Introduction}
For a graph $G$, let $w(G)$ be its whiskering and let $A(w(G))$ be the
Artinian squarefree monomial algebra associated with the independence complex
of $w(G)$. Holleben and Nicklasson conjectured that
\[
  \alpha(G)\geq3\quad\Longrightarrow\quad A(w(G))\text{ does not have the
  weak Lefschetz property}.
\]
They proved a range with large independence number and reported computations
for every graph on at most seven vertices and for many graphs on eight
vertices~\cite[Corollary~3.12, Conjecture~3.15, Remark~3.16]{Holleben_2026}.

We close the eight-vertex case over fields of characteristic zero. The field
qualification is explicit here because the ranks of the relevant integer
matrices can change in positive characteristic, while the source conjecture
does not state a characteristic.

\begin{theorem}\label{thm:main}
Let $K$ be a field of characteristic zero and let $G$ be a finite simple graph
on eight vertices. If $\alpha(G)\geq3$, then $A_K(w(G))$ does not have the
weak Lefschetz property.
\end{theorem}

The proof has two parts. Existing structural results, together with a short
Hilbert-function estimate, cover $\alpha(G)\geq4$. For $\alpha(G)=3$ we give
an exhaustive exact certificate: all 6,021 isomorphism classes have a
noninjective multiplication map $A_4\to A_5$.

\section{The algebra and its multiplication matrices}
Write $V(G)=[8]$. The graph $w(G)$ has vertices
$x_1,\ldots,x_8,y_1,\ldots,y_8$, the edges $x_ix_j$ for
$\{i,j\}\in E(G)$, and the eight whisker edges $x_iy_i$. Define
\[
A_K(w(G))=\frac{K[x_1,\ldots,x_8,y_1,\ldots,y_8]}
{(x_i^2,y_i^2,x_iy_i,\ x_ix_j:\{i,j\}\in E(G))}.
\]
For an independent set $X$ of $G$ and $Y\subseteq[8]\setminus X$, put
\[
 e_{X,Y}=\prod_{i\in X}x_i\prod_{j\in Y}y_j.
\]
The elements $e_{X,Y}$ with $|X|+|Y|=d$ form a basis of $A_d$. In
particular, if $i_k(G)$ is the number of independent $k$-sets of $G$, then
\begin{equation}\label{eq:hilbert}
 h_d:=\dim_K A_d=\sum_{k=0}^{\alpha(G)}i_k(G)\binom{8-k}{d-k}.
\end{equation}

Let $L=\sum_{i=1}^8(x_i+y_i)$ and denote the matrix of
$\times L:A_d\to A_{d+1}$ by $M_d(G)$. Its rows and columns are indexed by
the above monomial bases, and
\begin{equation}\label{eq:incidence}
 M_d(G)[(X',Y'),(X,Y)]=1
 \quad\Longleftrightarrow\quad
 \operatorname{supp}(e_{X,Y})\subset
 \operatorname{supp}(e_{X',Y'}),
\end{equation}
with support sizes $d$ and $d+1$; all other entries are zero.

For completeness, testing this one form is sufficient. Indeed, if a
Lefschetz form exists over an infinite field, the common maximal-rank locus
is a nonempty Zariski-open subset of $A_1$. It therefore meets the torus on
which every coefficient is nonzero. Diagonal rescaling of the variables is
an automorphism of this monomial quotient and takes any such form to $L$.
Thus $A$ has WLP if and only if every $M_d(G)$ has maximal rank.

\section{Structural reduction}
\begin{proposition}\label{prop:structural}
If $G$ has eight vertices and $\alpha(G)\geq4$, then $A_K(w(G))$ fails WLP
in characteristic zero.
\end{proposition}
\begin{proof}
If $\alpha(G)\geq5$, the assertion is Holleben--Nicklasson
Corollary~3.12~\cite{Holleben_2026}, since $\alpha(G)\geq 8/3+2$.

Suppose $\alpha(G)=4$. Their Theorem~3.10, applied to an independent
four-set, says that $M_5:A_5\to A_6$ is not
surjective~\cite[Theorem~3.10]{Holleben_2026}. Formula
\eqref{eq:hilbert} gives
\[
 h_6-h_5=-5i_2(G)+2i_4(G)-140.
\]
An independent four-set contains six independent pairs, so $i_2(G)\geq6$,
while $i_4(G)\leq\binom84=70$. Hence $h_6-h_5\leq-30$. Maximal rank
would therefore require $M_5$ to be surjective, which it is not.
\end{proof}

It remains to treat $\alpha(G)=3$. In this stratum, \eqref{eq:hilbert}
also gives $h_5-h_4=5i_2+5i_3-14>0$, so any nonzero vector in
$\ker M_4(G)$ proves failure of maximal rank. Indeed, an independent triple
gives $i_2\geq3$ and $i_3\geq1$.

\section{Exact finite certificate}
The program \texttt{geng 8} from nauty~2.7r3
\cite{McKay_2014,NautyGuide}
enumerated one graph6 representative of each of the 12,346 isomorphism
classes of simple graphs on eight vertices; the total agrees with the
published online census of McKay~\cite{McKayData}. Direct enumeration of all 256
vertex subsets gave the distribution in Table~\ref{tab:alpha}.
\begin{table}[ht]
\centering
\caption{Eight-vertex graph classes by independence number.}\label{tab:alpha}
\begin{tabular}{c@{\qquad}rrrrrrrr}
\toprule
$\alpha$&1&2&3&4&5&6&7&8\\
\midrule
classes&1&409&6021&4985&842&80&7&1\\
\bottomrule
\end{tabular}
\end{table}

For each of the 6,021 representatives with $\alpha=3$, the certificate file
stores the graph6 string, the ordered Hilbert vector, and a nonzero sparse
integer vector $v_G$ satisfying
\begin{equation}\label{eq:certificate}
 M_4(G)v_G=0.
\end{equation}
The verifier parses the graph6 string, recomputes $\alpha(G)$, reconstructs
the monomial bases and the inclusion relation \eqref{eq:incidence}, checks
the dimensions against the stored Hilbert vector, and evaluates every
integer coordinate of \eqref{eq:certificate}. It also rejects zero vectors,
repeated sparse indices, malformed dimensions, missing graph records, and
duplicate graph records. All 6,021 records pass.

\begin{proof}[Proof of Theorem~\ref{thm:main}]
Proposition~\ref{prop:structural} covers $\alpha(G)\geq4$. When
$\alpha(G)=3$, the exact certificate supplies a nonzero vector in
$\ker M_4(G)$ for every isomorphism class. Since $h_4<h_5$, the map is not
of maximal rank. Relabeling $G$ only permutes matrix rows and columns, so the
conclusion holds for every labelled graph. Finally, an integer matrix has
the same rank over every characteristic-zero field as over $\mathbb Q$.
\end{proof}

\section{Scope and data availability}
The theorem is a finite characteristic-zero result; it does not settle the
conjecture for graphs on nine or more vertices, nor does it assert a uniform
positive-characteristic statement. The certificate, verifier, and supporting
data are available in the supplementary archive distributed with this article.
\enlargethispage{7\baselineskip}
\bibliographystyle{amsplain}
\bibliography{references}
\end{document}
